\chapter{Data: sources, preparation and exploratory analysis}
\label{cap:datos}

\section{Data sources}

The main dataset of this work comes from real email marketing campaigns in the Spanish market, supplied in CSV format with a semicolon separator. The records arrived in two files, identified by two anonymised codes (LOW and CPC) that correspond to the two people who collected the data; they are not, therefore, different campaign types. For the analysis both files are unified into a single set, since they share the same structure and capture the same type of user--email interaction.

Each record captures an interaction event between a user and a campaign email. The information available per record covers user data (email address, date of birth, sex and postal code), data on the advertised product (sector, product name, brand and supplier company) and the type of event recorded, which can be \textit{click}, \textit{open} or \textit{ignored}. The \texttt{cpl} field records the cost per lead (CPL) of the advertised product and is only populated in one of the two sources (the one identified as CPC).

\section{Preprocessing and cleaning}

\subsection{Column selection and normalisation}

From among the fields available in the original dataset, twelve variables relevant to the analysis were selected. The columns were renamed for greater clarity and consistency: \texttt{fecha\_de\_nacimiento} was renamed \texttt{birthday}, \texttt{codigo\_postal} as \texttt{cp\_num}, \texttt{event\_timestamp} as \texttt{timestamp}, and likewise the remaining descriptive fields of the product and the user.

The text fields (\texttt{email}, \texttt{subject} and \texttt{cp\_num}) were normalised by removing spaces and converting to lower case. The categorical variables were stored as type \texttt{string} to avoid implicit conversions during processing.

\subsection{Filtering and anonymisation}

Records with a missing email address or postal code were discarded, as well as those with an age outside the 18--120 year range. To safeguard users' privacy, the identifiers were generated by means of an SHA-256 hash function: the \texttt{id\_user} field from the clean email address, and the \texttt{id\_product} field from the email subject. In this way referential integrity between tables is maintained without storing personally identifiable information.

\subsection{Event deduplication}

A given user may have interacted with the same email on multiple occasions. To avoid duplicates, deduplication was applied to the (\texttt{email}, \texttt{subject}) pairs, retaining the highest-priority event according to the hierarchy \textit{click} $>$ \textit{open} $>$ \textit{ignored}. This criterion reflects the informational relevance of each event type: a click implies active intent and has greater predictive value than a mere open.

\subsection{Normalisation of product categories}

The original values of the \texttt{product} field exhibited high textual variability, with multiple denominations for the same type of product. A normalisation function was implemented based on keyword scoring applied over the concatenation of \texttt{subject}, \texttt{product}, \texttt{sector} and \texttt{brand}, which maps each record to one of 25 standardised categories (including \texttt{other}). Among the resulting categories are \texttt{car\_insurance}, \texttt{health\_insurance}, \texttt{energy\_plan}, \texttt{mobile\_plan} and \texttt{funeral\_insurance}, among others. After normalisation, the catalogue was reduced to 70 unique products.

\section{Dimensional structure of the dataset}

The outcome of the cleaning process is three relational tables that constitute the dimensional model of the project.

The \textbf{users\_base} table holds one record per unique user, with their identification attributes and basic demographic data. The \textbf{products} table holds one record per unique campaign subject, with the descriptive attributes of the advertised product and its normalised category. The \textbf{events} table holds each interaction event, linking users and products through their respective identifiers and recording the event type and the moment in time at which it occurred. After the cleaning process, the events table contains 50.087 records corresponding to 37.815 unique users and 70 products, with a ratio of 11.940 clicks (23,8\%) against 38.147 opens (76,2\%). This ratio is the result of balanced sampling over the raw data, the justification for which is detailed in Section~\ref{sec:eda} (``Resampling comparison'').

\section{External enrichment}

The campaign dataset only records information on email address, date of birth, sex and postal code. With the aim of endowing each user with a more complete demographic and socio-economic profile, an external enrichment process was carried out using Spanish public data sources, mainly from the Instituto Nacional de Estadística (INE) and from municipal administrative registers.

\subsection{Geographic resolution}

The first step of the enrichment consisted of translating each user's postal code into their municipality of residence. To this end, a correspondence file was used covering the 11.692 postal codes currently in force in Spain. For the districts of Madrid and Barcelona, specific dictionaries were employed that allow resolution at district level, achieving greater geographic granularity in the two cities with the highest concentration of users. Of the total of 37.815 users, 93,6\% (35.405) were successfully linked to a specific municipality; the remaining 6,4\% (2.410 users), with no postal code\,$\rightarrow$\,municipality correspondence, were discarded from the enriched dataset. These 35.405 users with a complete geographic profile constitute the population on which the subsequent modules operate.

\subsection{Estimated demographic variables}

The additional demographic variables were generated by means of probabilistic sampling from municipal and provincial statistical distributions of the INE, so that each user receives values consistent with the demographic characteristics of their area of residence. This approach makes it possible to have complete profiles without the need for observed individual data, assuming that users are representative of the population of their municipality or province.

Employment status was estimated in two steps: first it was determined whether the user belongs to the active population on the basis of the municipal activity rate; for those active, the unemployment rate was applied to distinguish between employed and unemployed. The result was 53,3\% of users employed, 41,0\% inactive and 5,7\% unemployed. In an analogous manner, marital status, household type, the number of rooms of the dwelling and the probability of vehicle ownership were estimated from the corresponding provincial or municipal distributions.

\textbf{Fundamental limitation.} It is important to stress that these variables are \emph{ecological imputations}: they are not observed attributes of each individual, but values sampled from the aggregated distributions of their area of residence. Consequently, they share above all the \emph{geographic} signal of the user and do not reflect their real individual variability. Attributing to a specific user the statistical properties of their area constitutes a case of \emph{ecological fallacy}~\cite{robinson1950, greenland2001}, so the \emph{individual} predictive power of these variables is, by construction, limited. This limitation conditions the interpretation of the propensity results (Chapter~\ref{cap:propension}) and is revisited in the conclusions.

Table~\ref{tab:variables_enriquecidas} summarises the added variables, their source and the assignment method employed.

\begin{table}[htbp]
\centering
\caption{Variables added through external enrichment}
\label{tab:variables_enriquecidas}
\begin{tabular}{lll}
\toprule
\textbf{Variable} & \textbf{Source} & \textbf{Method} \\
\midrule
\texttt{labor\_status}   & INE, activity/unemployment rates 2025T4   & Bernoulli per municipality \\
\texttt{labor\_type}     & INE, municipal sectors          & Multinomial per municipality \\
\texttt{civil\_status}   & INE, provincial marital status       & Multinomial per province \\
\texttt{tiene\_coche}    & Municipal vehicle stock      & Bernoulli per municipality \\
\texttt{size\_hogar}     & INE, housing census            & Multinomial per municipality \\
\texttt{num\_room}       & INE, rooms per dwelling     & Multinomial per province \\
\bottomrule
\end{tabular}
\end{table}

\section{Final dataset}

The dataset resulting from preprocessing and external enrichment is composed of 35.405 users with a complete demographic profile of 16 variables, 70 products categorised into 25 families and 50.087 interaction events distributed over approximately twelve months (52 weeks). The user-product interaction matrix is very sparse (Section~\ref{sec:eda}, ``Sparsity and activity per user''), which rules out classical collaborative filtering methods in favour of approaches based on user characteristics and product content, as detailed in Chapter~\ref{cap:recomendador}.

\section{Exploratory data analysis}
\label{sec:eda}

Before modelling, an exploratory analysis was carried out to characterise the target variable, the user profiles, the behaviour by sector and the structure of the interactions. Its conclusions condition the subsequent methodological decisions.

\subsection{Target variable}

The variable to be predicted is binary: \textit{click} (active interaction) versus \textit{open} (without click). Of the 50.087 events, 11.940 are clicks (23,8\,\%) and 38.147 opens (76,2\,\%) (Figure~\ref{fig:eda_target}). This ratio corresponds to the \emph{balanced sample} used for training, not to the real \textit{click rate} of the complete dataset (see the following subsection, ``Resampling comparison''). Even so, the problem remains imbalanced, which justifies the use of PR-AUC as the main metric.

\begin{figure}[htbp]
\centering
\includegraphics[width=0.55\textwidth]{fig_eda_target}
\caption{Distribution of the target variable in the training dataset: clicks versus opens.}
\label{fig:eda_target}
\end{figure}

\subsection{Resampling comparison}

The \emph{raw} dataset exhibits an extreme imbalance: of the 2.026.166 original events, only 2,2\,\% are clicks (1 click for every $\approx$44 non-clicks). So that the propensity model has enough positive examples, a \textbf{1:1 balanced sampling} is applied (all clicks plus an equal number of non-clicks) which, after deduplication and the filters, leaves the training set at 50.087 events with 23,8\,\% clicks (Figure~\ref{fig:eda_resampling}). The bias thus introduced is corrected a posteriori with the \textit{prior}-shift correction (Chapter~\ref{cap:propension}), which returns the probabilities to the real production scale.

\begin{figure}[htbp]
\centering
\includegraphics[width=\textwidth]{fig_eda_resampling}
\caption{Response variable \textbf{before} (raw data, logarithmic scale) and \textbf{after} balanced sampling and cleaning. The resampling raises the click prevalence from 2,2\,\% to 23,8\,\%, endowing the model with enough positive examples.}
\label{fig:eda_resampling}
\end{figure}

\paragraph{Effect on the observed variables.} Balanced sampling
changes the prevalence of the response without distorting the composition of the sample.
Table~\ref{tab:antes_despues} and Figure~\ref{fig:eda_antes_despues} compare, before and after, the
\emph{observed} variables (not the synthetic demographic ones, which are added later). The structure is
preserved: \textbf{age} is practically identical (mean 49,2 vs.\ 49,6; median 49) and the
sector mix is maintained (insurance dominates with $\approx$40--46\,\% in both; the order of the
main sectors does not change). The only appreciable shift is in \textbf{gender}: the proportion
of men rises from 43,5\,\% to 49,1\,\%, an expected effect because clicks tend to come from men
(their \textit{CTR} is somewhat higher) and the sampling retains all clicks; the resampling
thus rebalances the response without materially biasing the profile of the sample.

\begin{table}[htbp]
\centering
\caption{Descriptive statistics of the observed variables: before (raw) vs.\ after resampling}
\label{tab:antes_despues}
\begin{tabular}{lrrrrr}
\toprule
\textbf{Set} & \textbf{Events} & \textbf{Click \%} & \textbf{Users} & \textbf{Mean age} & \textbf{Men \%} \\
\midrule
Before (raw)        & 2.026.166 & 2,2  & 306.476 & 49,2 & 43,5 \\
After (balanced) & 50.087    & 23,8 & 35.405  & 49,6 & 49,1 \\
\bottomrule
\end{tabular}
\end{table}

\begin{figure}[htbp]
\centering
\includegraphics[width=\textwidth]{fig_eda_antes_despues}
\caption{Observed variables before (grey) vs.\ after (red) resampling: age and gender per user, sector mix and response variable. The resampling rebalances the response while preserving age and sector; gender shifts slightly towards men.}
\label{fig:eda_antes_despues}
\end{figure}

\subsection{Demographic profile of the users}

Figure~\ref{fig:eda_demographics} summarises the main user variables. \textbf{Age} is concentrated between 35 and 64 years (mean 50, range 20--117); for modelling it is discretised into six bands (\texttt{age\_cat}: 18-24, 25-34, 35-44, 45-54, 55-64 and 65+), with 45-54 being the most populated. \textbf{Gender} is practically balanced. \textbf{Employment status} reflects the INE rates (around 53\,\% employed), \textbf{marital status} predominates in ``married'' and the \textbf{income class} (purchasing power index, IPA) is concentrated in the medium and medium-low levels. Apart from age and the geographic variables, these attributes are synthetic imputations, with the quality caveat already noted (``Fundamental limitation'').

\begin{figure}[htbp]
\centering
\includegraphics[width=0.92\textwidth]{fig_eda_demographics}
\caption{Demographic distributions of the users: age (continuous and categorised into the \texttt{age\_cat} bands of the M2 model), gender, employment status, marital status and income class (IPA).}
\label{fig:eda_demographics}
\end{figure}

\subsection{Behaviour by sector}

The click rate varies notably according to the \textbf{sector} of the product (Figure~\ref{fig:eda_ctr_sector}). \texttt{motor} is the sector with the highest propensity (0,37), followed by \texttt{servicios legales} (0,29) and \texttt{hogar} (0,27); \texttt{seguros}, the one with the largest volume, with 22.833 events, has a rate of 0,26, while \texttt{sorteos} (0,10) and \texttt{derechos humanos} (0,07) are the least effective. This heterogeneity between sectors suggested \emph{a priori} a hierarchical sector\,$\rightarrow$\,product architecture (first estimating the sector propensity and then refining the specific product). However, the comparison of the propensity module (Chapter~\ref{cap:propension}) shows that this decomposition worsens performance across all algorithms, so the sector is finally incorporated as one more variable of a flat model.

\begin{figure}[htbp]
\centering
\includegraphics[width=0.85\textwidth]{fig_eda_ctr_sector}
\caption{Click rate by sector (only sectors with $\geq$100 events). The dashed line marks the overall rate (0,238); in red, the sectors above the mean.}
\label{fig:eda_ctr_sector}
\end{figure}

\subsection{Click rate by demographic variables}

The \emph{real} demographic variables (age and gender, the only observed and non-imputed ones) barely discriminate the click propensity (Figure~\ref{fig:eda_ctr_demo}). By \textbf{gender}, men click somewhat more than women (0,26 versus 0,22). By \textbf{age band}, the rate ranges between 0,20 (25--44 years) and 0,30 (55--64 years), with a gentle U-shaped pattern; the 18--24 band, with an apparently high rate (0,31), barely has 179 events and is not conclusive. The demographic differences in the click rate are modest (of the order of $\pm0{,}05$ around the mean of 0,238) compared with the strong variation by sector (from 0,07 to 0,37). This contrast anticipates the central finding of the propensity module (Chapter~\ref{cap:propension}): the signal resides much more in the \emph{product} than in the user's \emph{profile}.

\begin{figure}[htbp]
\centering
\includegraphics[width=0.92\textwidth]{fig_eda_ctr_demo}
\caption{Click rate by age band (left) and by gender (right). The dashed line marks the overall rate (0,238). The demographic differences are modest compared with those of the sector.}
\label{fig:eda_ctr_demo}
\end{figure}

\subsection{Product catalogue and cost per lead (CPL)}

The catalogue is concentrated: the 25 categories are not distributed uniformly. Figure~\ref{fig:eda_products} (left) shows the twelve categories with the most events; \texttt{health\_insurance} (10.623 events), \texttt{energy\_plan} and \texttt{funeral\_insurance} lead the volume. This concentration is relevant for the recommendation: the three most clicked categories account for 45\,\% of all clicks and the top five for 60\,\%, which turns \emph{popularity} into a \textit{baseline} that is hard to beat (Chapter~\ref{cap:recomendador}). The colour of each bar indicates whether its CTR exceeds the mean: categories such as \texttt{vehicle\_purchase} (CTR 0,43) or \texttt{home\_security} (0,33) are very effective despite their lower volume.

The \textbf{CPL} (cost per lead, interpreted as the profit that each click yields in the ROI analysis) is available for 54 of the 70 products; in the records with no CPL reported it is imputed with the median of 8\,\euro. Its distribution (Figure~\ref{fig:eda_products}, right) is asymmetric, with most products between 5 and 12\,\euro\ and a tail up to 31\,\euro; the sectors with the highest mean CPL are \emph{salud y belleza} (19\,\euro) and \emph{hogar} (14,5\,\euro). These figures feed the economic analysis of Chapter~\ref{cap:negocio}.

\begin{figure}[htbp]
\centering
\includegraphics[width=\textwidth]{fig_eda_products}
\caption{Left: top 12 categories by event volume, coloured according to whether their CTR exceeds (red) or not (blue) the overall mean; the CTR of each one is annotated. Right: distribution of the CPL among the products that have it reported.}
\label{fig:eda_products}
\end{figure}

\subsection{Sparsity and activity per user}

The most important structural constraint of the project is the \textbf{scarcity of interactions per user} (Figure~\ref{fig:eda_sparsity}): 83,8\,\% of users have a single recorded event, the mean is 1,32 events and the maximum, 42. Combined with the 70 products, this produces a user--product matrix that is 98,1\,\% empty. The methodological consequence is direct: with so little individual history, classical collaborative filtering and a purely individual recommender are not viable, and approaches based on \emph{characteristics} (of the user and, above all, of the product) and on behavioural \emph{segments} are called for.

\begin{figure}[htbp]
\centering
\includegraphics[width=0.6\textwidth]{fig_eda_sparsity}
\caption{Number of events per user (grouping users with five or more into ``5+''). The vast majority interact only once.}
\label{fig:eda_sparsity}
\end{figure}

\subsection{Temporal evolution}

Activity is not uniform over time (Figure~\ref{fig:timeseries}): after a gradual start, the volume of clicks is concentrated between August and December 2025, preceded by a period of low activity (fewer than 50 clicks per week) before 28 July 2025. This concentration conditions the temporal forecasting module (M4), which only has around 22 weeks of active history.

\begin{figure}[htbp]
\centering
\includegraphics[width=\textwidth]{fig_timeseries}
\caption{Weekly evolution of the click volume (January 2025 -- January 2026). The grey shaded band delimits the inactive period prior to the start of the active campaign (28 July 2025).}
\label{fig:timeseries}
\end{figure}

\subsection{Data quality}

\paragraph{Cleaning funnel.} The final dataset is the result of a chain of successive filters over the raw data, summarised in Table~\ref{tab:funnel}. Each step responds to an explicit quality criterion (removing records without key information, safeguarding privacy, avoiding duplicates and ensuring a resolvable geographic profile). The most relevant drop in users occurs when requiring a postal code\,$\rightarrow$\,municipality correspondence (6,4\,\% not resolvable), discarded so as not to introduce arbitrarily imputed geographic profiles.

\begin{table}[htbp]
\centering
\caption{Cleaning funnel: from the raw data to the final modelling dataset}
\label{tab:funnel}
\begin{tabular}{p{6.2cm}rrr}
\toprule
\textbf{Stage} & \textbf{Events} & \textbf{Clicks} & \textbf{Users} \\
\midrule
Raw data (all raw files)                             & 2.026.166 & 44.997 (2,2\,\%) & --- \\
1:1 balanced sample + cleaning (dedup email/subject, age 18--120, missing) & 50.087 & 11.940 (23,8\,\%) & 37.815 \\
After requiring postal code\,$\rightarrow$\,municipality resolution    & 50.087 & 11.940 & 35.405 \\
\bottomrule
\end{tabular}
\end{table}

\paragraph{Completeness and consistency.} After cleaning, the tables present no missing values in the variables used by the models (the only unpopulated column is \texttt{district\_name}, not used in modelling); the CPL, absent in part of the records, is imputed with the median (8\,\euro) only for the ROI analysis. No duplicates were detected in the dimensional tables after the event deduplication. The coverage of the geographic enrichment reaches 100\,\% of the 35.405 final users. The main quality caveat is the one already noted: the \emph{fine-grained} demographic variables (employment, household, vehicle, etc.) are ecological imputations whose individual quality is bounded by construction (``Fundamental limitation'').
